\documentclass[11pt]{article}

% "final" prints the author block; "review" would anonymise it and add line numbers.
\usepackage[final]{acl}
\usepackage{times}
\usepackage{latexsym}
\usepackage[T1]{fontenc}
\usepackage[utf8]{inputenc}
\usepackage{microtype}
\usepackage{booktabs}
\usepackage{amsmath}
\usepackage{amssymb}
\usepackage{graphicx}
\usepackage{url}

% ============================================================
% TITLE AND AUTHORS
% ============================================================

\title{EcoExplain-VLM: Evaluating Vision--Language Models for Classifying
and Explaining Forest Vegetation Change in Kanha Tiger Reserve}

\author{
Rama Chandra Prasad Pillutla\thanks{Associate Professor. Email:
\texttt{rcprasad@iiit.ac.in}} \\
Lab for Spatial Informatics \\
International Institute of Information Technology Hyderabad, India
\And
Chinta Phani Rama Vaibhav\thanks{Email:
\texttt{chinta.vaibhav@students.iiit.ac.in}} \\
Lab for Spatial Informatics \\
International Institute of Information Technology Hyderabad, India
\And
Shaik Afzal Basha\thanks{Email:
\texttt{afzalbasha.shaik@students.iiit.ac.in}} \\
Lab for Spatial Informatics \\
International Institute of Information Technology Hyderabad, India
}

\begin{document}

\maketitle

% ============================================================
% ABSTRACT
% ============================================================

\begin{abstract}
Monitoring vegetation change in protected forests requires combining
evidence from several sensors, yet it is unclear whether general-purpose
vision--language models (VLMs) can classify such change and justify
their decisions. This study presents EcoExplain-VLM, a benchmark of 366
forest patches in Kanha Tiger Reserve, Madhya Pradesh, India. For each
patch, matched dry-season data (January--February 2022 and 2023) from
Sentinel-2, Sentinel-1 SAR, CHIRPS, ERA5-Land and NASADEM were encoded
as a pair of before/after RGB images and a structured evidence card.
Labels (stable, decline or increase) are percentile ranks of the
patch-mean NDVI change, and the patches were split into training,
development and test sets (256/55/55). Four open VLMs with 0.5--3
billion parameters were evaluated zero-shot on the same 55 test
patches. Qwen2-VL-2B reached 56.4\% accuracy, while LLaVA-Interleave-0.5B,
PaliGemma-3B and BLIP-2 OPT-2.7B scored 32.7--36.4\%, no better than
always predicting the majority class (34.5\%). A single NDVI-change
threshold fitted on the training set reaches 94.5\%. Low-rank adaptation
(LoRA) of Qwen2-VL-2B on 16 training patches for one epoch raised
accuracy to 67.3\%. Both Qwen2-VL-2B figures were obtained with the full
evidence card, whose final summary line contains the candidate label.
Explanations were then generated from cards with the label removed and
scored automatically by checking each claimed direction of change against
the measured value. Free-form prompts produced few checkable claims or
claims that conflicted with the evidence. A structured prompt that
restricts the model to citing supplied facts raised claim correctness to
75.8\% for Qwen2-VL-2B and 71.1\% for LLaVA-Interleave-0.5B, although
the justifications mostly restate those facts. A pilot LoRA on 16
filtered self-generated justifications reached 61.8\% class accuracy in
explanation mode but only 50.5\% correct claims. Small open VLMs
therefore do not yet extract the change signal that a one-number rule
captures, and they need constrained prompting to stay consistent with
the evidence.
\end{abstract}

\noindent
\textbf{Keywords:} vision--language models; forest change detection;
multi-source remote sensing; Sentinel-2; Sentinel-1 SAR; zero-shot
evaluation; LoRA; explanation faithfulness; Kanha Tiger Reserve

% ============================================================
% 1. INTRODUCTION
% ============================================================

\section{Introduction}

Forest managers in protected landscapes need to know not only
\emph{whether} a patch of canopy has become greener or browner between
two dates, but also \emph{why} that reading is credible. Optical indices
such as NDVI \citep{rouse1974,tucker1979}, radar backscatter, rainfall,
soil-moisture anomalies and terrain together constrain that
interpretation \citep{lillesand2015}. Operational change products
already exist at continental and global scale
\citep{hansen2013,kennedy2010,zhu2014}, but they do not produce a short
natural-language justification that a non-specialist can check against
the underlying numbers.

Vision--language models (VLMs) map images and text into a shared
representation and generate free text
\citep{liu2023llava,wang2024qwen2vl,beyer2024,li2023blip2}. In
principle, such a model could read a pair of dry-season image chips and
a compact evidence card, and answer with one of three labels
(\texttt{stable}, \texttt{decline}, \texttt{increase}) together with a
justification. In practice, small open models hallucinate content
\citep{li2023pope}, produce explanations that do not reflect how they
reached their answer \citep{turpin2023,jacovi2020}, and are sensitive to
prompt wording \citep{zhao2021calibrate}. Remote-sensing VLMs such as
GeoChat \citep{kuckreja2024} show that domain-specific training helps,
but they have not been evaluated on multi-sensor forest-change evidence
of the kind used here.

This study asks three questions for a single, compact reserve:
\begin{enumerate}
\item Can open VLMs with 0.5--3 billion parameters classify relative
vegetation change from before/after RGB chips and an evidence card
without task-specific training?
\item How does a small LoRA update \citep{hu2022lora} change that
behaviour?
\item With the label removed from the card, do the models' written
justifications agree with the measured changes, and does prompt
structure or fine-tuning on justifications improve that agreement?
\end{enumerate}

The contribution is a fully specified benchmark, EcoExplain-VLM, for
Kanha Tiger Reserve, together with an evaluation that reports its own
weaknesses. On the 55-patch test set, only Qwen2-VL-2B exceeds the
majority-class baseline (56.4\% against 34.5\%), and a 16-patch LoRA
raises it to 67.3\%. Both figures use an evidence card that ends with a
pipeline summary line containing the candidate label. The other three
models answer almost every patch with the same class, while a
one-number NDVI-change rule reaches 94.5\%. For explanations, a
structured ``cite-only'' prompt raises claim correctness to the 70\%
range, but the resulting texts restate the supplied facts rather than
interpret the images. Fine-tuning Qwen2-VL-2B on 16 of its own filtered
justifications gives 61.8\% class accuracy and 50.5\% correct claims.

The study does not produce an operational disturbance map, and it does
not claim that current small VLMs are ready for reserve management. The
labels rank NDVI change within one region and one pair of years, and a
human check of 19 patches agreed with them on only 8. These limits are
part of the result.

% ============================================================
% 2. DATA AND METHODS
% ============================================================

\section{Data and Methods}

\subsection{Study Area}

Kanha Tiger Reserve lies in the Maikal hills of central India, in the
Mandla and Balaghat districts of Madhya Pradesh. Its core zone covers
917.44\,km$^2$ and its buffer zone 1,134.36\,km$^2$, 2,051.79\,km$^2$ in
total \citep{kanha2026}. The forests are dominated by sal
(\emph{Shorea robusta}) and mixed deciduous stands, interspersed with
grassland meadows \citep{kanha2026}. Most rain falls in the
June--September monsoon, while January--February is cool, dry and
relatively cloud-free; both observation windows therefore use those
months.

Processing used the rectangle
$80.433^\circ$--$81.050^\circ$\,E,
$22.117^\circ$--$22.450^\circ$\,N (2,350.7\,km$^2$ in Earth Engine). The
rectangle is larger than the gazetted reserve, so some patches may lie
in the buffer zone or just outside the official boundary. All 1,200
candidate patch centres lie inside it. The candidate pool has an empty
strip near $80.75^\circ$--$80.85^\circ$\,E, inherited from the earlier
repository that created the pool; the reason was not recorded and is
most likely a forest-cover filter.

\subsection{Data Sources}

\subsubsection{Optical and Synthetic Aperture Radar Data}

Sentinel-2 Level-2A surface reflectance \citep{drusch2012} provides the
optical bands at 10--20\,m. Clouds and shadows were masked with the
Scene Classification Layer, and median composites were formed for
1~January--28~February 2022 (T0; 20 scenes) and the same window in 2023
(T1; 23 scenes). From each composite we computed NDVI, EVI
\citep{huete2002}, NDMI \citep{gao1996} and NBR \citep{keybenson2006},
and rendered a true-colour RGB chip for the VLMs.

Sentinel-1 GRD C-band backscatter \citep{torres2012} was composited for
the same windows from the dominant orbit direction (descending; 15
scenes at T0 and 20 at T1). The VV and VH channels (in dB) respond to
canopy structure and moisture and are unaffected by cloud. All exports
and composites were produced in Google Earth Engine \citep{gorelick2017}.

\subsubsection{Ancillary Environmental Data}

CHIRPS daily rainfall \citep{funk2015} was summed over each 58-day
window. ERA5-Land \citep{munozsabater2021} supplies 2\,m air temperature
and top-layer (0--7\,cm) volumetric soil moisture, expressed as
anomalies from a 2015--2021 climatology for the same months. NASADEM
\citep{nasadem2020} provides elevation, slope and aspect at 30\,m. None
of these layers enters the automatic label. They appear on the evidence
card so that a justification can mention drought or terrain when the
numbers support it.

\subsection{Preprocessing and Patch Generation}

Each patch is built by drawing a circle of 320\,m radius around a
candidate centre (half of $64 \times 10$\,m), taking its bounding
square, and clipping that square to the study rectangle, so patches on
the edge are slightly smaller. Exported GeoTIFFs are typically
$70 \times 65$ pixels in EPSG:4326 (pixel size $0.0000898^\circ$). Fifteen
rasters are stored per patch, 5,490 files in total. After cloud masking,
the masked fraction is 0.8--2.3\% per patch; candidates with a masked
fraction above 0.15 in either year were labelled uncertain and never
sampled. For the VLMs, the T0 and T1 RGB chips were resampled to
$512 \times 512$ pixels.

The nearest-neighbour distance between the 366 selected centres is at
least 639\,m (adjacent, non-overlapping squares), with a median of
1,430\,m. The data split is a random stratified draw rather than a
spatial block split \citep{ploton2020}, so spatial autocorrelation
between training and test patches remains a limitation.

\subsection{Change Evidence and Labelling}

Let $\Delta$ be the patch-mean NDVI change,
$\mathrm{NDVI}_{T1} - \mathrm{NDVI}_{T0}$. Over the 1,200-candidate
pool, the 15th and 85th percentiles of $\Delta$ are $q_{15} = -0.0377$
and $q_{85} = +0.0031$. The automatic label is
\begin{equation}
\mathrm{class} =
\begin{cases}
\text{uncertain} & \text{mask} > 0.15 \text{ in T0 or T1},\\
\text{decline} & \Delta \le q_{15},\\
\text{increase} & \Delta \ge q_{85},\\
\text{stable} & |\Delta| < 0.05,\\
\text{uncertain} & \text{otherwise}.
\end{cases}
\end{equation}
Because both percentiles lie within $\pm 0.05$, every unmasked patch
between $q_{15}$ and $q_{85}$ is labelled stable, and the final
``uncertain'' case does not occur for this pool. The region as a whole
lost a little greenness (mean $\Delta = -0.016$, s.d.\ $0.027$), so
``increase'' means the greenest 15\% of the pool rather than recovery in
an absolute sense: increase patches have $\Delta$ between $-0.001$ and
$+0.075$, and half of them changed by less than 0.011. ``Decline''
means the 15\% with the largest drop ($\Delta$ of about $-0.034$ to
$-0.071$). Only 34 of the 366 patches have $|\Delta| \ge 0.05$. The
classes are therefore relative ranks for this reserve and pair of
years.

The first 300 patches (100 per class) were drawn from the candidate
pool. A second batch added up to 25 patches per class, keeping new
centres at least 960\,m (1.5 patch widths) from existing ones; only 22
decline candidates passed this filter, giving 372 patches. Six patches
with incomplete exports were removed, leaving 366: 120 decline, 124
stable and 122 increase. The percentiles for the second batch were
recomputed on the remaining pool, so the class ranges overlap slightly
(decline patches reach $\Delta = -0.034$ and stable patches go down to
$-0.036$).

Each patch also has a structured evidence card listing index values and
changes, radar means, climate anomalies, terrain and a one-line pipeline
summary. That summary line includes the string
\texttt{candidate\_class=}\ldots, which is the field used as the gold
label. The card's qualitative trend words (for example
\texttt{moderate\_decline}) were computed from the $512 \times 512$
tensors with fixed thresholds, not from the percentile rule, and they
disagree with the label on 108 of 366 patches. The card's NDVI values
are on average 0.92 times the Earth Engine means. Bounding boxes from
Otsu thresholding exist on only three patches, and all stored narrative
explanations are template text, because no language-model explanation
was generated when the dataset was built.

A human reference class was assigned to 19 patches (13 training, 3
development, 3 test). It agrees with the automatic label on 8 of 19
(42\%; Wilson 95\% interval 23--64\%; \citealp{wilson1927}). An earlier
36-patch pilot agreed on 17 of 36 (47\%). There was no second
annotator. The automatic labels are therefore a reproducible ranking,
not a validated disturbance map.

\subsection{Dataset Construction and Data Splits}

Each patch is stored as one JSONL record containing image paths,
per-source statistics, the evidence card and the labels. The patches
were split 256/55/55 into training, development and test sets by
stratified sampling on the label (seed 42). The test set contains 18
decline, 18 stable and 19 increase patches, so always predicting
``increase'' gives a majority-class baseline of $19/55 = 34.5\%$. Every
result in this paper uses this test set.

\subsection{Vision--Language Model Evaluation}

Four instruction-tuned open models were compared:
Qwen2-VL-2B-Instruct \citep{wang2024qwen2vl},
LLaVA-Interleave-0.5B \citep{li2024interleave},
PaliGemma-3B-mix-224 \citep{beyer2024} and
BLIP-2 OPT-2.7B \citep{li2023blip2,zhang2022opt}, all loaded from
Hugging Face. Parameter counts are those of the released checkpoints.
Decoding is greedy. A deterministic parser maps the output to a class:
it reads a single-word answer directly and otherwise searches the text
for the first class word or synonym (for example, ``decrease'' is read
as decline). All 55 outputs were parsed for every model.

\subsubsection{Zero-Shot Evaluation}

The main classification setting, Condition~C, gives a model both RGB
chips and the evidence card and asks for one word: \texttt{decline},
\texttt{increase} or \texttt{stable}. The four models did not receive
identical input (Table~\ref{tab:inputs}). Qwen2-VL-2B and
LLaVA-Interleave-0.5B received both images and the full card, including
the summary line with \texttt{candidate\_class}. PaliGemma refused to
answer long analytical prompts, so it received both images, the card's
NDVI line and a short instruction. BLIP-2 accepts a single image, so the
two chips were placed side by side in one image with a short prompt and
no evidence text. The zero-shot runs were carried out on a laptop.

\begin{table}[t]
\centering
\small
\begin{tabular}{lll}
\toprule
Model & Images & Text \\
\midrule
Qwen2-VL-2B & T0, T1 & full card \\
LLaVA-Interleave-0.5B & T0, T1 & full card \\
PaliGemma-3B & T0, T1 & NDVI line only \\
BLIP-2 OPT-2.7B & stitched & none \\
\bottomrule
\end{tabular}
\caption{Input per model in Condition~C. ``Full card'' includes the
pipeline summary line containing \texttt{candidate\_class}.}
\label{tab:inputs}
\end{table}

Three supporting analyses help interpret the results. First, an
\texttt{evidence\_only} variant removes the summary line and any other
line that names the label. This changes the input only for LLaVA,
because the PaliGemma and BLIP-2 prompts never contained the line, and
Qwen2-VL-2B was run in this mode on only three patches. Second, a
closed-set probe scores LLaVA by the likelihood it assigns to each class
word, with a stated decision rule and a class bias fitted on the
development set. Third, two baselines that use no model apply a
threshold to the $\Delta$NDVI value on the card: one with fixed cut-offs
of $\pm 0.02$, and one with cut-offs fitted on the training set.

\subsubsection{Parameter-Efficient Fine-Tuning}

LoRA \citep{hu2022lora} was applied to Qwen2-VL-2B with rank 8,
$\alpha = 16$ and dropout 0.1 on the attention and MLP projection
layers, using AdamW \citep{loshchilov2019} with a learning rate of
$5 \times 10^{-5}$, weight decay 0.05 and gradient accumulation over
four steps. Training ran on a GTX 1080 Ti node of the institute's Ada
cluster. Two separate runs are reported.

\paragraph{Class LoRA.} The adapter was trained on 16 training patches
for one epoch with the classification target and the full Condition~C
card (final training loss 1.84), and tested on the full test set. This
is a small-scale configuration, not a complete fine-tune of the 256
training patches. Exploratory runs on all 256 training patches for two
to three epochs reached 100\% development and test accuracy, because
the model learned to copy \texttt{candidate\_class} from the summary
line. These runs are excluded from the results.

\paragraph{Explanation LoRA.} Qwen2-VL-2B itself acted as teacher: it
generated a justification for each training patch, and a generation was
kept only if its class matched the label, it contained at least two
checkable claims, and all of those claims agreed with the measured
changes. Of 256 training patches, 28 passed. A larger 7B teacher could
not be used because its download exceeded the cluster's disk quota. The
training run was capped at the first 128 training patches, so 16 of the
28 filtered justifications were used. Both training and evaluation used
the card without the label, the checkpoint was selected on the
development set, and the test set was evaluated in explanation mode.

\subsubsection{Evaluation Metrics}

Classification is scored by accuracy with Wilson 95\% intervals
\citep{wilson1927}, macro-F1, and comparison with the 34.5\%
majority-class baseline; McNemar's test \citep{mcnemar1947} compares
paired predictions.

Explanations are scored automatically. The scorer searches the text for
an indicator (NDVI, EVI, NBR, NDMI, rainfall, temperature, soil
moisture, VV, VH) and pairs it with the nearest direction word
(increase-type, decrease-type or stable-type). Each measured change is
converted to a direction using an indicator-specific dead-band, below
which a change counts as stable. A claim is \emph{correct} when its
direction matches the measured direction. We report the mean number of
checkable claims per explanation, the share of correct claims (pooled
over patches), the share of explanations with at least one incorrect
claim, and self-consistency (whether the dominant claimed direction
agrees with the model's own class). ROUGE-L \citep{lin2004rouge} is
computed but not used as a measure of faithfulness, because overlap with
the supplied text rewards copying rather than reasoning
\citep{jacovi2020}. The scorer measures consistency with the evidence,
not faithfulness to the model's internal reasoning.

\subsection{Experimental Workflow}

\begin{enumerate}
\item Export the multi-sensor stacks and RGB chips from Earth Engine.
\item Assign automatic labels, build evidence cards and freeze the
256/55/55 split.
\item Run zero-shot Condition~C classification for the four models.
\item Run the \texttt{evidence\_only}, closed-set and $\Delta$NDVI
baseline analyses.
\item Train and test the class LoRA (16 patches) on Qwen2-VL-2B.
\item Generate explanations from label-free cards with a free-form
prompt, then with a structured cite-only prompt, and score them.
\item Filter teacher justifications and train and test the explanation
LoRA.
\end{enumerate}

The development set was used only for prompt wording, the closed-set
bias and the explanation LoRA checkpoint; no headline result was chosen
by looking at the test set.

% ============================================================
% 3. RESULTS AND DISCUSSION
% ============================================================

\section{Results and Discussion}

\subsection{Dataset Characteristics}

Table~\ref{tab:data} summarises the benchmark. The optical index changes
move together: $\Delta$NDVI correlates with $\Delta$NBR, $\Delta$NDMI
and $\Delta$EVI at 0.80, 0.77 and 0.78. Mean window rainfall fell from
69.1\,mm in 2022 to 25.1\,mm in 2023, the temperature anomaly rose from
$-1.52^\circ$C to $+0.88^\circ$C, and the soil-moisture anomaly fell
from $+0.024$ to $-0.043$. Mean VV and VH backscatter changed by $-0.47$
and $-0.48$\,dB. On average, therefore, the second dry season was drier
and slightly less green, which is why ``increase'' is a relative rank
and why many increase patches have $\Delta$ close to zero.

\begin{table}[t]
\centering
\small
\begin{tabular}{lr}
\toprule
Quantity & Value \\
\midrule
Study rectangle & 2,350.7\,km$^2$ \\
Candidate pool / final patches & 1,200 / 366 \\
Decline / stable / increase & 120 / 124 / 122 \\
Train / development / test & 256 / 55 / 55 \\
Test decline / stable / increase & 18 / 18 / 19 \\
Majority baseline (test) & 34.5\% \\
Human agreement & 8/19 (42\%) \\
Mean $\Delta$NDVI (s.d.) & $-0.016$ (0.027) \\
Patches with $|\Delta\mathrm{NDVI}| \ge 0.05$ & 34 \\
Card trend words disagree with label & 108 \\
\bottomrule
\end{tabular}
\caption{EcoExplain-VLM dataset summary. Labels are NDVI-change ranks
within this region and pair of years.}
\label{tab:data}
\end{table}

\subsection{Zero-Shot Model Performance}

Table~\ref{tab:zeroshot} reports zero-shot accuracy on the test set.
Qwen2-VL-2B is the only model above the majority baseline, at 56.4\%
(31/55; 95\% interval 43.3--68.6\%). Even this margin is not
statistically secure at $n = 55$: McNemar's test against always
predicting ``increase'' gives $p = 0.12$. The other three models do not
separate the classes. LLaVA-Interleave-0.5B answered ``decrease'' for
all 55 patches, PaliGemma answered ``decline'' for 42 of 55, and BLIP-2
repeated its own prompt, from which the parser read the first class word
(``decline''). BLIP-2's 32.7\% is therefore a parsing artefact rather
than a prediction.

\begin{table*}[t]
\centering
\small
\begin{tabular}{llrcr}
\toprule
Model or baseline & Input & Acc. & 95\% CI & Macro-F1 \\
\midrule
Qwen2-VL-2B & images + full card & \textbf{56.4\%} & 43.3--68.6 & 0.466 \\
PaliGemma-3B & images + NDVI line & 36.4\% & 24.9--49.6 & 0.271 \\
LLaVA-Interleave-0.5B & images + full card & 32.7\% & 21.8--45.9 & 0.164 \\
BLIP-2 OPT-2.7B & stitched image only & 32.7\% & 21.8--45.9 & 0.164 \\
\midrule
LLaVA-Interleave-0.5B, \texttt{evidence\_only} & images + card without label line & 36.4\% & 24.9--49.6 & 0.251 \\
LLaVA-Interleave-0.5B, closed-set probe & class-word likelihood, dev-fitted bias & 60.0\% & 46.8--71.9 & 0.479 \\
\midrule
Majority class (always ``increase'') & --- & 34.5\% & 23.4--47.7 & 0.171 \\
$\Delta$NDVI rule, $\pm 0.02$ cut-offs & card $\Delta$NDVI & 63.6\% & 50.4--75.1 & 0.606 \\
$\Delta$NDVI rule, cut-offs fitted on train & card $\Delta$NDVI & 94.5\% & 85.1--98.1 & 0.945 \\
\bottomrule
\end{tabular}
\caption{Classification on the test set ($n = 55$). The top block is
generative zero-shot Condition~C. The middle block contains supporting
runs that are not zero-shot generative results. The bottom block uses no
model.}
\label{tab:zeroshot}
\end{table*}

Two observations follow. First, the result is not explained by model
size alone: Qwen2-VL-2B and LLaVA-Interleave-0.5B received the same
input, and only the larger of the two used it. Model size cannot be
compared with PaliGemma or BLIP-2 because their inputs were shorter.
Second, the full-card condition is not a pure test of visual or spectral
understanding. Qwen2-VL-2B may be reading the summary line, the index
values, the images or a mixture of these, and it has not yet been tested
on the full test set without the summary line. For LLaVA, removing the
line made almost no difference (32.7\% to 36.4\%), because it did not
use the card in either condition.

The baselines show that the card contains enough information to solve
the task: a single threshold on its $\Delta$NDVI value reaches 94.5\%
with cut-offs fitted on the training set. This is expected, since the
labels are defined from NDVI change, and it sets a clear target that
none of the VLMs approaches. The closed-set probe lifts LLaVA to 60.0\%,
which suggests that constraining the output helps a small model even
when free generation fails, although the probe never predicts
``stable''.

\subsection{Effect of Fine-Tuning}

The class LoRA raises Qwen2-VL-2B from 56.4\% to \textbf{67.3\%} (37/55;
95\% interval 54.1--78.2\%; Table~\ref{tab:lora}). The intervals of the
two runs overlap, so with 16 training patches and $n = 55$ test patches
this is an indicative gain rather than a powered comparison. The
fine-tuned model was correct on all 18 decline and all 19 increase test
patches, but it labelled every stable patch as decline. The adapter
therefore learned two of the three classes and lost the third.

\begin{table}[t]
\centering
\small
\begin{tabular}{lcc}
\toprule
Qwen2-VL-2B run & Train & Test acc. \\
\midrule
Zero-shot (full card) & 0 & 56.4\% \\
Class LoRA, 1 epoch (full card) & 16 & 67.3\% \\
\midrule
\multicolumn{3}{l}{\emph{Excluded or exploratory}} \\
Full card, 2--3 epochs & 256 & 100\% \\
Trained with card, tested on images only & 128 & 32.7\% \\
\bottomrule
\end{tabular}
\caption{Class LoRA results for Qwen2-VL-2B. The 100\% runs copied the
label from the summary line and are not counted as results.}
\label{tab:lora}
\end{table}

Two exploratory runs support this reading of the summary line. Adapters
trained on all 256 patches with the full card reached 100\% because
they could read the answer from the card. An adapter trained with the
card on 128 patches and then tested on images alone scored 32.7\%,
below the majority baseline. Neither run is a valid estimate of
performance, and the second mixes training and test conditions, but
together they suggest that the fine-tuned model relies on the card text
rather than on the images. The explanation LoRA is a
different task and is discussed in Section~\ref{sec:expl}.

\subsection{Error Analysis}

Zero-shot Qwen2-VL-2B was correct on all 18 decline patches and on 13 of
18 stable patches, but it answered ``stable'' for all 19 increase
patches. For about half of the increase class the card's trend words
say ``stable'', and the model followed those words rather than the
summary line. Errors therefore concentrate where the label is a weak
rank: increase patches with $\Delta$ near zero are reasonably called
stable, and decline and stable patches overlap slightly after the
second-batch percentile shift.

The 108 cards whose trend words disagree with the label build a
contradiction into the input: a model that reads the descriptive text
faithfully can still miss the gold label. Together with the low human
agreement (8 of 19), this means that some ``errors'' are likely to be
label noise rather than model failure. The test set was not relabelled
after seeing model outputs.

The failures of the other models are mostly failures to follow the task.
LLaVA always produced the same answer, PaliGemma refused long prompts
and largely answered with one class, and BLIP-2 echoed its instruction.
The parser accepted the echo because it searched the whole output for a
class word; a stricter parser should read only newly generated text and
reject copied instructions.

\subsection{Explanation Quality}
\label{sec:expl}

All explanation runs used the card without the label line.
Table~\ref{tab:explain} summarises them. With the first, free-form
prompt, which asked the model to describe and justify the change in its
own words, PaliGemma answered with a class or ``unanswerable'' and no
reason, and BLIP-2 repeated the prompt, so neither produced a checkable
claim. LLaVA wrote fluent but repetitive text in which 56.7\% of the
checkable claims were correct. Qwen2-VL-2B was run on only three
patches in this mode, and one of its claims misstated the rainfall
change. Free-form output from these models is not usable as an audit
trail.

\begin{table*}[t]
\centering
\small
\begin{tabular}{lrrrrr}
\toprule
Setting and model & Class acc. & Claims / expl. & Correct claims & $\ge 1$ wrong & Self-consistent \\
\midrule
Free-form, LLaVA-Interleave-0.5B & 29.1\% & 1.09 & 56.7\% & --- & --- \\
Free-form, PaliGemma-3B & 33.3\%$^{a}$ & 0 & --- & --- & --- \\
Free-form, BLIP-2 OPT-2.7B & 32.7\%$^{b}$ & 0 & --- & --- & --- \\
\midrule
Structured, Qwen2-VL-2B & 63.6\% & 2.85 & \textbf{75.8\%} & 36.4\% & 100\% \\
Structured, LLaVA-Interleave-0.5B & 45.1\%$^{c}$ & 2.33 & \textbf{71.1\%} & 34.5\% & 90.7\% \\
\midrule
Explanation LoRA, Qwen2-VL-2B & 61.8\% & 6.84 & 50.5\% & 74.5\% & 96.4\% \\
\bottomrule
\end{tabular}
\caption{Explanation results on the test set, label removed from the
card. $^{a}$12 of 36 answered patches. $^{b}$Prompt echo; every answer
parsed as ``stable''. $^{c}$23 of 51 answered patches. The free-form
Qwen2-VL-2B run (3 patches, 43.8\% correct claims) is omitted.}
\label{tab:explain}
\end{table*}

The structured prompt gives the model a numbered list of allowed facts
that must be copied verbatim as supporting evidence, a line stating an
NDVI decision rule for the patch ($\pm 0.02$), and a request for a
one-sentence justification about NDVI only. Under this prompt, claim
correctness rises to 75.8\% for Qwen2-VL-2B and 71.1\% for
LLaVA-Interleave-0.5B. Most of the remaining incorrect claims are not
invented facts: 35 of Qwen's 38 and 35 of LLaVA's 37 flagged claims
describe a copied change smaller than the scorer's dead-band (for
example, an EVI change of $-0.006$ described as a decrease). High claim
correctness here means that the model did not invent numbers; it does
not mean that the model reasoned about the change
\citep{jacovi2020,turpin2023}. Qwen2-VL-2B's class matched the stated
rule on 53 of 55 patches, and the rule alone scores 63.6\%, exactly
Qwen's accuracy in this mode. This is instruction-following, not a new
classification ability.

The explanation LoRA reached 61.8\% class accuracy (34/55; 95\% interval
48.6--73.5\%) and produced more claims per explanation (6.84), but only
50.5\% of them were correct and 74.5\% of explanations contained at
least one incorrect claim. This is worse than structured zero-shot
prompting. The bottleneck is the teacher: only 28 of 256 of its own
generations passed the filter, and only 16 were used. With this
evidence, self-distillation of a 2B model does not teach faithful
explanation; the adapter mainly learned the output format.

\subsection{Comparison with Existing Approaches}

Pixel- and segment-based change detectors
\citep{hansen2013,kennedy2010,zhu2014,daudt2018} optimise map accuracy
rather than a three-class patch rank with a sentence, and they remain
the right tool for a wall-to-wall disturbance map. Remote-sensing
question answering \citep{lobry2020} and change captioning
\citep{liu2022levircc} pair imagery with language, but they work on
single-sensor optical imagery and do not check generated text against
measured values. EcoExplain-VLM asks a different question: whether a
general VLM can use the same multi-sensor evidence that a human analyst
would see, and whether its explanation agrees with that evidence.

Remote-sensing VLMs such as GeoChat \citep{kuckreja2024} are trained on
overhead imagery and spatial referring expressions and are a natural
next baseline. They were not evaluated here because the available
compute limited the study to general checkpoints of 0.5--3 billion
parameters. CLIP-style encoders \citep{radford2021} could score the
image pair without generating text, but they would not produce the
justification that the benchmark is designed to test.

The most informative comparison this dataset supports is with simple
baselines. Against the majority class, only Qwen2-VL-2B with the full
card shows a clear margin. Against a single $\Delta$NDVI threshold, none
of the models comes close.

% ============================================================
% 4. LIMITATIONS AND FUTURE DIRECTIONS
% ============================================================

\section{Limitations and Future Directions}

\textbf{Labels.} The classes are percentile ranks of one-year
dry-season NDVI change, not validated disturbance, degradation or
recovery, and the absolute changes are small. Human agreement on 19
patches is 42\%, with a single annotator. An expert audit with a second
annotator is needed before any management use.

\textbf{Label in the input.} The Condition~C card contains
\texttt{candidate\_class}, and both headline Qwen2-VL-2B results (56.4\%
and 67.3\%) use it. Future evaluations should use the
\texttt{evidence\_only} card by default, which the code now supports.

\textbf{Unequal inputs.} PaliGemma and BLIP-2 received shorter input
than Qwen2-VL-2B and LLaVA, so the four models are not compared under
identical conditions.

\textbf{Sample size and split.} With $n = 55$, 95\% intervals are about
$\pm 13$ percentage points wide. The random stratified split does not
prevent spatial leakage \citep{ploton2020}; a spatial block split should
be added as a second test set.

\textbf{Models.} The results cover four small open VLMs and say nothing
about larger proprietary or remote-sensing-pretrained models.

\textbf{Explanation scoring.} Automatic claim matching can miss
paraphrases, can pair a direction word with the wrong indicator, and can
mark a true but small change as incorrect because of the dead-band. The
reported claim correctness is not a human judgement of faithfulness.

\textbf{Geography.} One reserve, one pair of years and one season were
studied. Transfer to other parks, wet-season composites or fire years is
unknown.

Five extensions would change the scientific picture rather than add
rows to a leaderboard: (i)~make label-free evaluation the only primary
table, including Qwen2-VL-2B on all 55 test patches; (ii)~add a spatial
block test; (iii)~collect expert labels for all 366 patches or a
sufficiently large subsample; (iv)~fine-tune on all 256 training
patches without the label line, with early stopping on the development
set; and (v)~evaluate a larger or remote-sensing-pretrained VLM, and a
stronger teacher for explanation distillation, on the same cards.

% ============================================================
% 5. CONCLUSIONS
% ============================================================

\section{Conclusions}

EcoExplain-VLM packages 366 Kanha patches, matched 2022/2023 dry-season
Sentinel-2 and Sentinel-1 data, climate and terrain context, and a
three-class NDVI-change label into a reproducible VLM benchmark. On the
55-patch test set, Qwen2-VL-2B reaches 56.4\% zero-shot accuracy with
the full evidence card, and a 16-patch class LoRA raises this to
67.3\%; both figures depend on a card that contains the candidate label.
LLaVA-Interleave-0.5B, PaliGemma-3B and BLIP-2 OPT-2.7B stay at
32.7--36.4\%, the level of always guessing one class, while a single
$\Delta$NDVI threshold reaches 94.5\%. For explanations, free-form
prompts give unusable output, and a structured cite-only prompt makes
explanations consistent with the evidence (75.8\% and 71.1\% correct
claims) at the cost of restating the supplied facts. An explanation LoRA
trained on 16 self-generated justifications gives 61.8\% class accuracy
and 50.5\% correct claims and is a pilot rather than an improvement.

Small open VLMs, as evaluated here, are not a substitute for spectral
change detection in Kanha. Under a constrained prompt they can repeat
the facts they are given. The value of the benchmark is that it makes
this limitation measurable and gives larger or domain-adapted models a
fixed, label-free test against which to show real progress.

% ============================================================
% ACKNOWLEDGEMENTS
% ============================================================

\section*{Acknowledgements}

The authors thank the Lab for Spatial Informatics, International
Institute of Information Technology Hyderabad, and the institute's Ada
computing cluster, which was used for LoRA training. Sentinel-1 and
Sentinel-2 data are provided by the European Space Agency through the
Copernicus programme; CHIRPS, ERA5-Land and NASADEM are used under their
respective open licences. Kanha Tiger Reserve is administered by the
Madhya Pradesh Forest Department and the National Tiger Conservation
Authority. This study uses satellite and reanalysis data only and is not
an official assessment of the reserve.

% ============================================================
% REFERENCES
% ============================================================

\bibliography{references}

% ============================================================
% APPENDIX
% ============================================================

\clearpage
\appendix

\section{Additional Experimental Details}

\subsection{Fine-Tuning Configuration}

Table~\ref{tab:ftconfig} lists the settings shared by the class and
explanation LoRA runs. Both used Qwen/Qwen2-VL-2B-Instruct with the
Hugging Face \texttt{transformers} and \texttt{peft} libraries, and both
trained for one epoch on 16 examples. For the class LoRA, the 16
patches were the configured training size. For the explanation LoRA,
the 16 were the filtered justifications that fell within the first 128
training patches.

\begin{table}[h]
\centering
\small
\begin{tabular}{ll}
\toprule
Setting & Value \\
\midrule
Base model & Qwen2-VL-2B-Instruct \\
LoRA rank / $\alpha$ / dropout & 8 / 16 / 0.1 \\
Target modules & q, k, v, o, gate, up, down \\
Optimiser & AdamW \\
Learning rate / weight decay & $5\times10^{-5}$ / 0.05 \\
Gradient accumulation & 4 steps \\
Epochs / training examples & 1 / 16 \\
Hardware & GTX 1080 Ti (Ada cluster) \\
\bottomrule
\end{tabular}
\caption{LoRA configuration.}
\label{tab:ftconfig}
\end{table}

\subsection{Prompt and Output Parsing}

\paragraph{Classification.} Qwen2-VL-2B and LLaVA-Interleave-0.5B
received the two images, the full evidence card and the instruction
``Reply with exactly one word after the images: decline or increase or
stable. Do not copy the evidence text.'' PaliGemma received the two
images, the card's NDVI line and ``Two forest RGB chips: first image T0,
second T1. Vegetation change class --- answer with exactly one word:
stable, decline, or increase.'' BLIP-2 received the stitched image and
``Side by side: left forest patch T0, right forest patch T1. One word
only --- decline, increase, or stable:''. The \texttt{evidence\_only}
variant removes every card line that names \texttt{candidate\_class}.

\paragraph{Explanations.} The free-form prompt tells the model that
image 1 shows the patch in January--February 2022 and image 2 in
January--February 2023, asks it to describe the visible change in one
sentence, decide whether vegetation is stable, declined or increased,
and justify the decision. Qwen2-VL-2B and LLaVA answer in JSON;
PaliGemma and BLIP-2 receive a one-sentence version of the request. The structured prompt adds a
numbered list of allowed facts to be copied verbatim, a line stating
the $\pm 0.02$ NDVI rule for the patch, and a one-sentence justification
limited to NDVI. An automatic check confirms that no explanation prompt
contains \texttt{candidate\_class}.

\paragraph{Parsing.} The parser lowercases the output, accepts a
single-word answer directly, and otherwise searches for the first class
word or synonym. It does not exclude text copied from the prompt, which
explains BLIP-2's echo artefact.

\paragraph{Teacher filtering.} For the explanation LoRA, a generation
was kept only if its parsed class equalled the label, it contained at
least two checkable claims, and every claim agreed with the measured
change. This filter kept 28 of 256 training patches.

\subsection{Additional Results}

With $n = 55$, intervals are wide: 31/55 (56.4\%) corresponds to
43.3--68.6\% and 37/55 (67.3\%) to 54.1--78.2\%. Per-class counts
(about 18 patches per class) are too small to support a detailed error
taxonomy. The closed-set LLaVA probe (60.0\%) relies on a
development-fitted bias and is reported as a probe, not as a zero-shot
result; the same probe left PaliGemma and BLIP-2 at 32.7\%, predicting
``stable'' for every patch. In the \texttt{evidence\_only} runs, LLaVA
answered ``decline'' for 51 patches and ``increase'' for 4.

\end{document}
